% section: 具体的な型から特殊関数へ

\section{具体的な型から特殊関数へ}

\subsection{「型を解く」ことの先にあるもの}

前半の実戦編では、式の形を見て変数を置き換えた。
定数項があれば不動点を引き、べきがあれば対数を取り、
三角関数型なら倍角公式を使った。

ここで、もう一歩だけ先へ進む。
その型が一つの問題に現れただけなら、解法を知っていれば十分である。
しかし、同じ型が波、確率、近似、幾何の中で何度も現れるなら、
その解法の中心にある関数を独立した数学的対象として研究したくなる。

\[
  \begin{array}{c}
    \text{具体的な漸化式}                      \\
    \downarrow\quad\text{置換・公式・帰納法}     \\
    \text{標準的な解}                        \\
    \downarrow\quad\text{別の問題にも繰り返し現れる} \\
    \text{特殊関数}
  \end{array}
\]

この流れを、いくつかの例で追う。

\subsection{階乗型の漸化式を実数へ延長する}

階乗は
\[
  (n+1)!=(n+1)n!
\]
を満たす。
この式の $n$ を実数 $x$ に広げたいなら、
\[
  F(x+1)=xF(x),\qquad F(1)=1
\]
を満たす関数を探すことになる。

この条件だけでは関数は一意に決まらない。
例えば、適切な周期関数を掛けても漸化式を保つからである。
そこで、正値性や滑らかさ、成長の仕方などの自然な条件を追加する。

代表的な選択が、
\[
  \Gamma(x)=\int_0^\infty t^{x-1}e^{-t}\,dt
\]
である。
部分積分により、
\[
  \begin{aligned}
    \Gamma(x+1)
     & =\int_0^\infty t^xe^{-t}\,dt      \\
     & =x\int_0^\infty t^{x-1}e^{-t}\,dt \\
     & =x\Gamma(x)
  \end{aligned}
\]
となる。
さらに $\Gamma(1)=1$ なので、
\[
  \Gamma(n+1)=n!
\]
である。

ここでは、実戦編の「階比型で積を整理する」という発想が、
整数列の外へ延長されている。
階乗に似た数列を解くための道具だった積分表示が、
ガンマ関数という特殊関数になった。

\subsection{三角関数型からチェビシェフ多項式へ}

実戦編で扱った
\[
  a_{n+1}=2\cos\theta\,a_n-a_{n-1}
\]
に対して、
\[
  a_n=\cos(n\theta)
\]
を倍角公式と帰納法で示した。

ここで $\theta$ を毎回書く代わりに、
\[
  x=\cos\theta
\]
と置き、
\[
  T_n(x)=\cos(n\theta)
\]
と定義する。
すると、
\[
  T_{n+1}(x)=2xT_n(x)-T_{n-1}(x)
\]
となる。

つまり、三角関数型の漸化式を、$x$ の多項式の漸化式として読み直した。
これがチェビシェフ多項式である。

この変換が有効なのは、三角関数が多項式近似の問題に姿を変えるからである。
例えば、$[-1,1]$ 上で関数を近似したいとき、
\[
  1,x,x^2,\ldots
\]
をそのまま使う代わりに、
\[
  T_0(x),T_1(x),T_2(x),\ldots
\]
を使うと、区間上の振動を制御しやすい。

\subsection{同じ三項漸化式から、別の多項式族が生まれる}

チェビシェフ多項式だけが特別なのではない。
係数の異なる三項漸化式から、別の自然な多項式族が生まれる。

例えばルジャンドル多項式は、
\[
  (n+1)P_{n+1}(x)
  =(2n+1)xP_n(x)-nP_{n-1}(x)
\]
を満たす。
母関数は
\[
  \frac1{\sqrt{1-2xt+t^2}}
  =\sum_{n=0}^{\infty}P_n(x)t^n
\]
である。

この多項式は、
\[
  (1-x^2)y''-2xy'+n(n+1)y=0
\]
の多項式解でもある。
さらに、
\[
  \int_{-1}^{1}P_m(x)P_n(x)\,dx=0
  \qquad(m\neq n)
\]
という直交性を持つ。

一つの漸化式を解くという入口から、
\[
  \text{母関数}
  \longleftrightarrow
  \text{微分方程式}
  \longleftrightarrow
  \text{直交性}
\]
という三つの見方が得られる。

\subsection{ガウス関数の微分からエルミート多項式へ}

正規分布に現れる
\[
  e^{-x^2}
\]
を何度も微分すると、
必ず多項式と $e^{-x^2}$ の積になる。
そこで、
\[
  H_n(x)e^{-x^2}
  =(-1)^n\frac{d^n}{dx^n}e^{-x^2}
\]
で $H_n$ を定める。

生成関数は
\[
  e^{2xt-t^2}
  =\sum_{n=0}^{\infty}H_n(x)\frac{t^n}{n!}
\]
であり、$t$ で微分して係数を比較すると、
\[
  H_{n+1}(x)=2xH_n(x)-2nH_{n-1}(x)
\]
が得られる。

ここでは「母関数を作る」という発展パターンが、
直交多項式の具体的な漸化式を作り出している。
ルジャンドル多項式では区間 $[-1,1]$、
エルミート多項式では全実数と重み $e^{-x^2}$ が現れる。
問題の対称性や重みが変わると、自然な特殊関数も変わる。

\subsection{円筒の波からベッセル関数へ}

円筒座標で波動方程式を解くと、
\[
  x^2y''+xy'+(x^2-\nu^2)y=0
\]
が現れる。
その解であるベッセル関数は、整数次数について
\[
  J_{n+1}(x)
  =\frac{2n}{x}J_n(x)-J_{n-1}(x)
\]
を満たす。

三角関数型や直交多項式型と同じく、
隣り合う次数の関数を結ぶ漸化式がある。
ただし、ここでは解は多項式ではなく、円筒の波を表す関数になる。

生成関数
\[
  \exp\left(\frac{x}{2}\left(t-\frac1t\right)\right)
  =\sum_{n=-\infty}^{\infty}J_n(x)t^n
\]
を $t$ で微分して係数比較を行うと、
\[
  \frac{x}{2}\{J_{n-1}(x)+J_{n+1}(x)\}
  =nJ_n(x)
\]
が得られ、上の漸化式が戻る。

ここで、実戦編の「隣の項を結ぶ漸化式を利用する」という発想が、
波動方程式の解を効率よく計算する道具になっている。

\subsection{具体的な型と特殊関数の対応}

\begin{center}
  \begin{tabular}{c|c|c}
    実戦での見抜き方 & 発展した対象    & 広がる理論     \\ \hline
    階比型・積の整理 & ガンマ関数     & 積分・確率     \\
    三角関数型    & チェビシェフ多項式 & 近似・直交性    \\
    三項漸化式    & ルジャンドル多項式 & 球面対称性     \\
    ガウス関数の微分 & エルミート多項式  & 正規分布・量子力学 \\
    円筒の波の漸化式 & ベッセル関数    & 波動・境界値問題  \\
  \end{tabular}
\end{center}

この表で重要なのは、関数の名前を増やすことではない。
目の前の漸化式に対して行った
\[
  \text{置く}\quad
  \text{変形する}\quad
  \text{生成する}\quad
  \text{直交化する}
\]
という操作が、別の問題でも繰り返し現れた結果として、
特殊関数が生まれたということである。
